% ==================== 六、问题四 ====================
\section{问题四：实时波动电价下的调度}

\subsection{问题分析与总体思路}

问题四指出外网电价也是\textbf{实时波动}的，要求在此条件下重算问题二（记为 4-2）
与问题三（记为 4-3）。与前三问的关系如下：

\begin{center}
\begin{tabular}{@{}l l l@{}}
\toprule
 & 电价 & 决策时点 \\
\midrule
问题二  & 每日固定（附件 1）    & 仅 $0$:00 \\
\textbf{4-2} & \textbf{实时波动（附件 4）} & 仅 $0$:00 \\
问题三  & 每日固定              & $0/6/12/18$ 时 \\
\textbf{4-3} & \textbf{实时波动}      & $0/6/12/18$ 时 \\
\bottomrule
\end{tabular}
\end{center}

因此 4-2 是“问题二 $+$ 电价随机性”，4-3 是“问题三 $+$ 电价随机性”。
两问均需解决一个新问题：\textbf{电价本身也未知，如何构造决策所需的场景？}

\subsection{4-2：波动电价下的日前调度}

\subsubsection{电价信息的假设（本问的关键判断）}

题面说明电价实时波动，但\textbf{未明确每天 $0$:00 是否已知全天价格}。
这个信息假设对结果影响很大，故我们保留两种情形并\textbf{分别完整求解}：

\begin{table}[H]
\centering
\caption{4-2 的两种电价信息情形}
\begin{tabular}{@{}l l l@{}}
\toprule
情形 & 假设 & 性质 \\
\midrule
\textbf{主模型} & 当天电价\textbf{未知}，只用此前日期预测 & 保守、符合“实时波动”的通常理解 \\
对照         & 当天 144 段电价已于 $0$:00 公布 & 明确放宽信息限制 \\
\bottomrule
\end{tabular}
\end{table}

主模型不把当天或未来的实际电价、负载、光伏作为制定当天计划的输入；
附件 4 的实际电价仅用于\textbf{事后结算}。

\subsubsection{电价预测器选择}

只用 1 月 8--31 日的 24 个历史预测日，比较四种规则：

\begin{enumerate}[itemsep=2pt,topsep=3pt]
  \item 前一日曲线；
  \item 近 7 日均值；
  \item 近 7 日形状 $+$ 前一日平均水平；
  \item 最近 4 次同星期几均值。
\end{enumerate}

以 1 月预测的 RMSE 选定，\textbf{2 月起固定方法不再更换}，避免用全年费用反向挑选参数。
入选方法为\textbf{最近 4 次同星期几均值}：

\begin{equation}
\hat p_{d,t}=\frac{1}{|\mathcal W_d|}\sum_{i\in\mathcal W_d}p_{i,t},
\qquad \mathcal W_d=\{d-7j:\ j=1,\dots,4,\ d-7j\ge0\}.
\end{equation}

预测 $d+1$ 日时进一步要求所有样本日期 $i<d$，不得使用尚未发生的第 $d$ 日价格。

\subsubsection{电价—净负载联合场景}

\textbf{核心论证：期望的不可分解性。}
当电价随机时，
\begin{equation}
\E\bigl[p_t(N_t-z_t)_+\bigr]\ \neq\ \E[p_t]\cdot\E\bigl[(N_t-z_t)_+\bigr].
\label{eq:q4-nonsep}
\end{equation}
不等式成立的原因是：\textbf{高电价可能与缺电同时发生}
（例如极端天气既推高需求又推高价格）。
若按式 \eqref{eq:q4-nonsep} 右端建模（即“平均电价近似”），
会系统性地\textbf{低估}缺电风险的实际代价。

因此场景\textbf{取自同一历史日}的两种残差配对：

\begin{equation}
N_{s,t}=\hat N_{d,t}+r^N_{s,t},\qquad
p_{s,t}=\max\bigl\{0.001,\ \hat p_{d,t}+r^p_{s,t}\bigr\},
\label{eq:q4-scen}
\end{equation}

其中 $r^N_{s,t}=N_{s,t}-\hat N_{s,t}$、$r^p_{s,t}=p_{s,t}-\hat p_{s,t}$ 均由历史日 $s$ 产生。
$0.001$ 元/kWh 是预先给定的正数下限，用于防止加性误差生成负电价，
被截断的场景数逐日记录以便审查。

场景集为 $\mathcal H_d=\{\max(7,d-56),\dots,d-1\}$，等权 $\omega_s=1/|\mathcal H_d|$。

\subsubsection{目标函数}

当天计划 $g,c,d$ 在所有场景间共享，紧急购电 $u_{s,t}$ 可随场景变化：

\begin{equation}
\min\ J_{42}=\sum_{t=1}^{144}\overline p_{d,t}\,g_t
+5\sum_s\omega_s\sum_{t=1}^{144}p_{s,t}\,u_{s,t}
+\underbrace{\sum_{t=1}^{144}\hat p_{d+1|d,t}\,g_{144+t}}_{\text{次日前瞻价值（不进报表）}},
\label{eq:q4-obj}
\end{equation}

其中 $\overline p_{d,t}=\sum_s\omega_s p_{s,t}$ 为场景平均电价，用于第一阶段购电费（计划量照付）。

\textbf{风险覆盖的分位解释}：固定某段的 $c,d$，令 $z=g+d-c$，
则其覆盖水平可用\textbf{电价加权分位}理解
\begin{equation}
\widetilde\omega_{s,t}=\frac{\omega_s p_{s,t}}{\overline p_{d,t}},\qquad
z\approx Q^{\widetilde\omega}_{0.8}(N_{s,t}),
\end{equation}
即高电价的场景会获得更高权重，从而抬高所需的供给水平。
但该式只是启发式解释，实际须与储能跨时段约束一并求解，不能逐段套用。

\subsubsection{约束与对照模型}

约束沿用问题二：计划下界 \eqref{eq:q2-floor}、场景兜底 \eqref{eq:q2-scen}、
SOC 与容量功率 \eqref{eq:q2-soc}--\eqref{eq:q2-cap}，以及 48 h 窗口终端条件。
另外设置\textbf{平均电价对照}：保留完全相同的净负载场景，
但把所有 $p_{s,t}$ 替换为 $\overline p_{d,t}$，以检验“配对关系”的实际作用。

\subsubsection{结果}

\begin{table}[H]
\centering
\caption{4-2 结果对照（2 月 1 日--12 月 31 日）}
\label{tab:q42}
\begin{tabular}{@{}l r r r@{}}
\toprule
方案 & {总费用（元）} & {较 Q2 重结算（元）} & {节省比例} \\
\midrule
原 Q2 计划按附件 4 重结算 & 15184704.23 & — & — \\
平均电价模型              & 15160921.45 & $-23782.79$ & $0.157\%$ \\
\textbf{联合场景模型（主）} & \bfseries 15163043.89 & \bfseries $-21660.34$ & $0.143\%$ \\
当日电价已知（对照）       & 15033043.18 & $-151661.05$ & $0.999\%$ \\
\bottomrule
\end{tabular}
\end{table}

\textbf{诚实说明}：在同年回测中，平均电价近似\textbf{略优于}联合场景
（差约 $2\,122$ 元）。我们不因此改用平均电价模型，理由有二：
(i) 联合场景保留了正确的概率结构（式 \eqref{eq:q4-nonsep} 说明它的期望是准确的），
(ii) 单一年度的 $2\,122$ 元差异不具统计显著性，而结构性正确性可跨年成立。
这一现象在结果文件中如实保留。

\textbf{“当日电价已知”的增益}：$151\,661.05$ 元（约 $1.0\%$），
说明\textbf{电价预测误差}是本问剩余成本的主要来源之一。

\begin{figure}[htbp]
\centering
\includegraphics[width=0.95\textwidth]{07_q42_price_information.pdf}
\caption{相对原第二问策略按附件 4 电价重新结算的费用节省。左图将节省分为
计划费与紧急费，菱形表示合计，负值表示该项费用增加；
右图展示均价预测与联合场景的月度净节省。当日电价已知为信息对照。}
\label{fig:q42-info}
\end{figure}

图 \ref{fig:q42-info} 显示月度节省存在正负变化：年度改善不保证各月均改善，
且均价预测方案在本年度略优（见上文说明）。该图同时说明
“信息更充分”与“当月费用更低”不总是同步。

\begin{figure}[htbp]
\centering
\includegraphics[width=0.95\textwidth]{08_q42_schedule_shift.pdf}
\caption{四个指定日的实际电价、优化采用的场景均价及购电计划差异。
左列阴影仅表示两条价格曲线的差距；右列为联合场景计划购电功率
减去原 Q2 计划购电功率，蓝色为增购、橙色为少购。}
\label{fig:q42-shift}
\end{figure}

将原第二问策略作为参照，可以观察到价格信息模型引入后购电时段发生重排。
场景均价与实际电价并不处处一致，价格曲线形状接近也不保证决策完全相同。
联合场景的决策还依赖价格与净负载误差的配对、储能约束及前期形成的状态，
因此左、右两列用于对照现象，不将每个购电差值单独归因于同一时段的价格预测误差。
此处展示的是日前计划，不应描述为当天观察实际价格后的日内修正。

\begin{figure}[H]
\centering
\includegraphics[width=0.9\textwidth]{q42_prices.pdf}
\caption{4-2：指定日期的实际电价曲线与预测电价}
\label{fig:q42-prices}
\end{figure}

\begin{figure}[H]
\centering
\includegraphics[width=0.9\textwidth]{q42_cost.pdf}
\caption{4-2：各方案费用对照}
\label{fig:q42-cost}
\end{figure}

\FloatBarrier

\subsection{4-3：波动电价 $+$ 多时点调整}

\subsubsection{模型构造}

4-3 在 4-2 的基础上叠加问题三的多时点调整机制：每天 $0$ 点定计划，
$6/12/18$ 点用新发布的光伏预报与更新的电价预测重新优化尚未执行的调度。
计费规则（调增 $1.5$ 倍、调减净成本 $-0.5$ 倍、紧急 $5$ 倍）与问题三完全一致，
即沿用式 \eqref{eq:q3-fee}。

\textbf{电价场景的处理}：每次发布时刻对未来电价采用
\textbf{此前 7 天同一时段均值}，并与同一历史日的净负载残差\textbf{配对}构造价格场景
（保留式 \eqref{eq:q4-scen} 的相关性结构）。次日采用历史价格预测。

\textbf{场景树、融合与执行}：与问题三相同——使用经验场景树保证非前视，
用在线融合整合历史预测与发布预报，只执行当前发布到下一更新前的控制量。

\subsubsection{结果}

\begin{table}[H]
\centering
\caption{4-3 主方案费用分解}
\label{tab:q43}
\begin{tabular}{@{}l r@{}}
\toprule
项目 & {金额（元）} \\
\midrule
原计划购电费       & 13457230.68 \\
增购费（$1.5$ 倍）   & 152250.38 \\
取消原费退款       & $-173200.59$ \\
取消违约费（$50\%$） & 86600.30 \\
净调整费（小计）     & 65650.08 \\
紧急购电费（$5$ 倍） & 1538854.67 \\
\midrule
\textbf{总费用} & \bfseries 15061735.43 \\
\bottomrule
\end{tabular}
\end{table}

\subsubsection{关键对比：调整机制价值的同口径分解}

固定电价（问题三）与波动电价（4-3）两问都建立了\emph{逐级增量}的对照方案，
故可对两问做\textbf{口径完全对称}的分解。下表左半为各方案总费用（元），
右半为由此得到的收益分解（万元）：

\begin{table}[H]
\centering
\small
\caption{两问逐级增量对照与同口径分解（结果期均 334 天）}
\label{tab:q4-compare}
\setlength{\tabcolsep}{4pt}
\begin{tabular}{@{}l r r@{\hskip 8pt}l r r@{}}
\toprule
方案 & 固定电价 & 波动电价 & 收益来源 & 固定 & 波动 \\
\midrule
无融合，仅 $0$ 点（B）    & 14389892.50 & 15516984.54
  & ① 在线预报融合 & $20.84$ & $22.19$ \\
融合预报，仅 $0$ 点       & 14181445.03 & 15295120.10
  & ② 多时点调整 & $15.90$ & \bfseries $23.34$ \\
融合 $+$ $0/6$ 点         & 14101371.47 & 15173094.46
  & \quad 加 $6$ 点  & $8.01$ & $12.20$ \\
融合 $+$ $0/6/12$ 点      & 14022358.10 & 15068610.47
  & \quad 加 $12$ 点 & $7.90$ & $10.45$ \\
\textbf{融合 $+$ $0/6/12/18$ 点} & \bfseries 14022396.98 & \bfseries 15061735.43
  & \quad 加 $18$ 点 & $-0.004$ & $+0.69$ \\
\bottomrule
\end{tabular}
\end{table}

“无融合 B 方案 $\to$ 主方案”的总改善可拆成\textbf{融合预报}与\textbf{多时点调整}
两部分，两者之和恒等于总改善（固定电价 $36.75$ 万、波动电价 $45.52$ 万）。
可见\textbf{在完全同口径下，波动电价显著放大了多时点调整的价值}：
由 $15.90$ 万元增至 $23.34$ 万元（\textbf{高约 $47\%$}），
而融合预报的价值仅由 $20.84$ 万元增至 $22.19$ 万元（高约 $7\%$）。
调整机制的相对重要性由 $43\%$（$15.90/36.75$）提升到 $51\%$（$23.34/45.52$），
\emph{超过了融合预报}。

这一点还可从\textbf{边际递减形态}上看出差异：固定电价下第三次调整
（加 $18$ 点）已经\textbf{零收益}（$-38.88$ 元），
而波动电价下仍有 $+6\,875$ 元正收益——说明电价波动\emph{推迟了}调整机制的饱和点。
原因在于电价波动引入了\textbf{第二个不确定性来源}：多时点更新不仅能修正光伏预报，
还能修正\textbf{电价预期}，当新信息显示“傍晚电价将冲高”时，
可及时调整购电与储能部署，这是固定电价下不存在的收益渠道。

需要强调：$4\text{-}2 \to 4\text{-}3$ 的\textbf{主方案}差值为 $101\,308.46$ 元，
但该差值横跨两问，同时包含“预测源切换”与“引入调整机制”两种效应，
\textbf{不能}作为调整机制的边际价值；上表的逐级增量才是同口径度量。
两问主方案的绝对费用仍相差约 $104$ 万元，主要来自附件 4 的实时电价整体高于附件 1。

\begin{figure}[htbp]
\centering
\includegraphics[width=0.95\textwidth]{09_q43_incremental_comparison.pdf}
\caption{4-3 与第三问主方案的费用差额及更新收益比较。左图按原计划费、
净调整费和紧急费分解 4-3 减第三问的总费用差；右图分别在两问内部
计算相邻更新方案之间的费用节省。}
\label{fig:q43-incremental}
\end{figure}

4-3 主方案费用为 $15\,061\,735.43$ 元，比第三问高 $1\,039\,338.45$ 元，
其中紧急费差额为 $877\,139.21$ 元，占总差额的 $84.39\%$。
该结果是不同电价输入及对应滚动策略共同形成的账面差异，
不能全部解释为电价波动或预测误差的因果成本。在各自内部的逐级对照中，
增加 6 点更新的节省由第三问的 $8.01$ 万元变为 4-3 的 $12.20$ 万元，
增加 12 点由 $7.90$ 万元变为 $10.45$ 万元；4-3 增加 18 点节省 $6\,875.04$ 元，
而第三问费用增加 $38.88$ 元。这表明新增调整机会的回测价值随运行情形而变化。

\begin{figure}[htbp]
\centering
\includegraphics[width=0.95\textwidth]{10_q43_schedule_differences.pdf}
\caption{四个指定日中，4-3 相对第三问的电价、最终执行购电功率及段末储电量差值。
蓝色表示 4-3 更高，橙色表示更低，白色表示接近零；各面板使用独立的对称色标，
数值未裁剪，竖虚线为 6、12、18 点。}
\label{fig:q43-diff}
\end{figure}

配对差值图保留了两问在相同实际负载与光伏条件下的时段对应关系。
电价、购电与储能状态的差异具有不同的时间分布，显示调度响应包含跨时段的能量转移。
各指定日的初始储电量已由此前滚动执行决定，两问并不相同，
因此状态差异既包含当日决策差异，也包含前期累积影响；
图中比较用于描述两套已保存运行结果，不将某一列的颜色直接解释为另一列的因果效应。

\begin{figure}[H]
\centering
\includegraphics[width=0.9\textwidth]{q43_cost.pdf}
\caption{4-3：各方案费用对照}
\label{fig:q43-cost}
\end{figure}

\begin{figure}[H]
\centering
\includegraphics[width=0.9\textwidth]{q43_dispatch.pdf}
\caption{4-3：指定日期的购电与储能调度（波动电价下）}
\label{fig:q43-dispatch}
\end{figure}

\FloatBarrier
